\begin{frame}[shrink=40]{장을 넘나드는 동치성: GDA vs 로지스틱회귀}
  \begin{columns}[T]
    \begin{column}{0.42\textwidth}
      \small
      \begin{tabular}{@{}lc@{}}
        \toprule
        측정 & 값 \\
        \midrule
        로지스틱 test AUC & 0.9934 \\
        GDA/LDA test AUC & 0.9855 \\
        출력 확률 상관계수 & 0.900 \\
        $\cos(w_{\text{logit}}, w_{\text{GDA}})$ & 0.203 \\
        test 라벨 일치율 & 93.6\% \\
        \bottomrule
      \end{tabular}
    \end{column}
    \begin{column}{0.55\textwidth}
      \centering
      \includegraphics[width=\linewidth]{ch08_3_gda_logistic_bars.png}
    \end{column}
  \end{columns}
  \vskip0.4em
  \textbf{결정 경계는 같고, 경계 벡터 방향은 다르다} --- \textbf{정규분포
  전제에 따라 동치성 범위가 달라진다}
\note{\begin{itemize}
\item 확인 문제 1의 답은 ``둘 다 $P(y{=}1\mid x)=\sigma(w^\top x+b)$
  라는 같은 \textbf{형태}에 도달한다''는 것이었다. ``동치''는
  어디까지 진실인가? breast\_cancer(569$\times$30)로
  \textbf{실측}(train/test 70/30, stratify, 스케일러
  train으로만):
  \small
\item \kb{결정 경계(분류)는 거의 같고}(라벨 93.6\% 일치, AUC 0.99
  부근), \kb{출력 확률도 꽤 비슷}하지만(상관 0.900),
  \kb{경계 벡터 $w$의 방향은 사실 다르다}($\cos 0.203$).
\item Week 3.2의 동치성은 \textbf{$P(x\mid y)$가 정규분포라는
  전제 하에} 성립하는 등식 --- 이 데이터가 정규분포가 아니므로
  전제가 \textbf{약하게 깨지고}, ``같은 \textbf{형태}의 경계,
  \textbf{다른 계수}''가 나온다.
\item {\footnotesize \kq{시험의 3점 답안: 정규분포 가정이 성립하면
  계수까지 같은 방향의 경계, 전제가 깨지면 경계의 \textbf{형태}
  (시그모이드-선형)만 같고 계수는 경로(생성적 vs 판별적)에 따라
  다르다 --- 전제·반례 한 쌍이 있는 답.} ``동치''를 ``계수도
  같다''로 지나치게 확장하지 않는 것이 정직한 답안.}
\end{itemize}}
\end{frame}
